\section{Résumé - français}

\textbf{Contexte :} Le pilotage de drones FPV en environnements complexes nécessite un 
retour vidéo en temps réel. Or, les fréquences standards (2.4 et 5.8 GHz) pénètrent mal 
les obstacles physiques (conditions NLOS). 

\textbf{Problématique :} La norme Wi-Fi HaLow (IEEE 802.11ah, Sub-1 GHz) offre une excellente 
pénétration, mais ses mécanismes natifs de fiabilité (acquittements MAC, retransmissions) 
induisent une latence et une gigue incompatibles avec les exigences strictes du vol en immersion.

\textbf{Méthode :} Ce travail implémente une architecture de bout en bout forçant une transmission 
dédiée au temps réel. Un pont matériel (SPI vers microcontrôleur STM32) encapsule un flux vidéo 
(optimisé en \textit{Bare-Metal} via MJPEG) en UDP Unicast. Cette approche est couplée à un forçage 
matériel de la Qualité de Service (QoS) pour écraser les délais d'accès au canal, verrouillant ainsi le 
modem radio sur une modulation robuste à basse latence.

\textbf{Résultat :} Après avoir isolé un temps de vol radio pur de l'ordre de 12.5 ms, l'approche itérative 
a permis d'éliminer les lourds goulots d'étranglement inhérents à l'encodage vidéo traditionnel. 
La latence visuelle globale absolue (\textit{Glass-to-Glass}) du système final a été mesurée à 
\textbf{40-50 ms en moyenne}. Ces performances valident définitivement la viabilité 
du Wi-Fi HaLow pour l'exploration par drone en milieux confinés.

\section{Abstract - English}

\textbf{Context :} FPV drone piloting in complex environments requires a real-time video feed. However, 
standard frequencies (2.4 and 5.8 GHz) struggle to penetrate physical obstacles (NLOS conditions).

\textbf{Problem :} The Wi-Fi HaLow standard (IEEE 802.11ah, Sub-1 GHz) offers excellent obstacle 
penetration, but its native reliability mechanisms (MAC acknowledgments, retransmissions) introduce 
latency and jitter incompatible with the strict requirements of immersive flight.

\textbf{Method :} This work implements an end-to-end architecture tailored for absolute real-time 
transmission. A hardware bridge (SPI to STM32 microcontroller) encapsulates a \textit{Bare-Metal} 
optimized MJPEG video stream into UDP Unicast. This approach is combined with a hardware Quality of 
Service (QoS) override to bypass channel access delays, locking the radio modem onto a robust, 
low latency modulation scheme.

\textbf{Result :} After isolating a pure radio transport time of approximately 12.5 ms, the iterative 
approach successfully eliminated the heavy bottlenecks inherent to traditional video encoding. 
The absolute overall \textit{Glass-to-Glass} visual latency of the final system was reduced to 
\textbf{an average of 40-50 ms}. These performances definitively validate the viability 
of Wi-Fi HaLow for drone exploration in confined environments.